\section{Introduction}
Small sacroiliac joint (SIJ) movements provide a way to examine how the pelvic ring transmits load. Articular geometry, ligamentous restraint and muscular compression contribute to stability, while force distribution determines how the ring is engaged \citep{Vleeming2012,Snijders1993,Hammer2013}. Mobility therefore cannot be specified without defining the loading context.

Pelvic form varies at several anatomical levels: overall scale, inlet and midpelvic proportions, pubic arch geometry and local SIJ morphology. Developmental variation includes sex differences in obstetrically relevant dimensions \citep{Huseynov2016}. At the joint itself, three-dimensional analysis of Japanese skeletal specimens demonstrated sex differences in iliac auricular shape \citep{Nishi2020}. These population contrasts encompass multiple anatomical features rather than a single mechanical property.

Size and shape must also be distinguished. Pelvic shape covaries with stature \citep{Fischer2015}, but sex differences in overall proportions and in specific features do not have the same relationship to size. Fischer and Mitteroecker found largely allometric dimorphism in pelvic height-to-width proportions and iliac orientation, whereas subpubic-angle dimorphism was largely non-allometric \citep{Fischer2017}. Sex, size and specific architecture are therefore not interchangeable descriptions of pelvic form.

The functional-anatomy question is how these levels of variation map onto mechanical response. Earlier finite-element comparisons of representative female and male models demonstrated sex-associated SIJ differences under standing loads \citep{Joukar2018}. A population cohort permits a complementary question: how do those contrasts change when dimensions are represented explicitly, and how does response vary with pelvic size under standardized force? Comparing adjusted coefficients describes associations without identifying causal pathways. This distinction also constrains functional interpretations of the relationship between pelvic support and obstetric accommodation \citep{Pavlicev2020}.

Individual architecture can be examined through matched geometry--density comparisons. Previous population modeling related SIJ auricular morphology to its mechanical environment \citep{HenysHammer2025}. Retaining individual geometry while standardizing the bone-density field tests whether geometry preserves subject-specific joint-motion outputs; retaining individual density on a common geometry supplies the reciprocal reference. These comparisons test simplification of represented bone density with soft-tissue properties held fixed, rather than separating all anatomical and material influences.

Repeated loading of the same subjects establishes the context in which scale and architecture are expressed. Bilateral and unilateral acetabular loads redistribute a fixed total force, while localized force pairs probe different sites on the ring. The latter are motivated by obstetric anatomy but do not simulate childbirth, which additionally involves fetal contact, soft tissues and changing boundary conditions \citep{ChenGrimm2021}. Using 278 imaging-derived pelves from the BoneDat workflow \citep{HenyssKuchar2025BoneDat}, we aimed to characterize SIJ response across anatomical scale, explicit architecture and load path. We examined conditional size--motion scaling under standardized prescribed forces, attenuation of sex-associated differences by explicit dimensions, and preservation of individual response under standardized density, within matched loading comparisons. The questions and directional hypotheses were retrospective and not preregistered; the scaling analysis retains exploratory statistical status.
